%%%%%%%%%%%%%%%%%%%%%%%%%%%%%%%%%%%%%%%%%
% MSc Environmental Engineering thesis poster
% Muhammad Faseeh Memon, supervised by Professor Michael Templeton
%
% Built to the frozen specification in FINAL_POSTER_BUILD_PACKAGE_25_AUG_2026.
% Asymmetric A1 landscape: left knowledge rail + wide results field.
%
% Built on the official Imperial College London Poster class (v1.0, 2024).
% The class file and both supplied sample templates are UNMODIFIED.
%
%!TEX program = xelatex
% Must be compiled with XeLaTeX (custom OpenType fonts).
%
% Scientific authority: GROUND_TRUTH_FINAL_SUBMITTED_THESIS.pdf, 293 pages,
% SHA-256 f371a8fe86ab6380bc5ee0426f9de708593277cb36d9bc6746b5cb14529df7aa.
% Every number traces to a page of that document; see poster_source_traceability.md.
%
% Panel widths deliberately sum to slightly LESS than the enclosing \linewidth.
% Summing to exactly \linewidth makes TeX wrap the second minipage onto a new
% line, which silently turns the poster into three pages.
%%%%%%%%%%%%%%%%%%%%%%%%%%%%%%%%%%%%%%%%%

\documentclass[
	landscape,   % Landscape page orientation
	a1papersize, % A1 (841 x 594 mm) rather than the class default A0
	print,       % Convert colours to CMYK and use the CMYK logo for printing
]{ImperialPoster}

% The supplied Imperial Sans Text family has no italic file, so the class's
% \setmainfont leaves the italic shape undefined and LaTeX substitutes a
% non-brand face. Redeclared here, in the poster file only, with a modest
% synthetic slant so that \textit stays in the brand typeface.
\setmainfont{ImperialSansText}[
	Path=Fonts/,
	UprightFont=*-Regular.ttf,
	BoldFont=*-Bold.ttf,
	ItalicFont=*-Regular.ttf,
	ItalicFeatures={FakeSlant=0.18},
	BoldItalicFont=*-Bold.ttf,
	BoldItalicFeatures={FakeSlant=0.18},
]

% The class A1 caption format is 10pt, below the official small-text level.
\DeclareCaptionFormat{postercaption}{\fontsize{16pt}{19pt}\selectfont#1#2#3}
\captionsetup{format=postercaption}

\definecolor{RuleGrey}{HTML}{D7DCE4}

%----------------------------------------------------------------------------------------
%	TYPE SCALE (05_TYPOGRAPHY_COLOUR_SPACING.md)
%----------------------------------------------------------------------------------------

\newcommand{\bodyfont}{\fontsize{22pt}{25pt}\selectfont}
\newcommand{\capfont}{\fontsize{16pt}{19pt}\selectfont}
\newcommand{\srcfont}{\fontsize{13.5pt}{15.5pt}\selectfont}
\newcommand{\reffont}{\fontsize{14pt}{15.5pt}\selectfont}

\newcommand{\railhead}[1]{%
	{\ImperialSansSemiBold\fontsize{30pt}{33pt}\selectfont\color{ICLDarkBlue}#1\par}%
	\vspace{1.2mm}{\color{ICLBlue}\rule{\linewidth}{1.5pt}}\par\vspace{1.5mm}}

\newcommand{\mainhead}[1]{%
	{\ImperialSansSemiBold\fontsize{36pt}{39pt}\selectfont\color{ICLDarkBlue}#1\par}%
	\vspace{1.2mm}{\color{ICLBlue}\rule{\linewidth}{2pt}}\par}

\newcommand{\blockhead}[1]{%
	{\bfseries\fontsize{26pt}{29pt}\selectfont\color{ICLDarkBlue}#1\par}\vspace{1.8mm}}

\newcommand{\callout}[1]{%
	{\ImperialSansExtraBold\fontsize{34pt}{36pt}\selectfont\color{ICLDarkBlue}#1\par}}

\newcommand{\srctag}[1]{\vspace{1.4mm}{\srcfont\color{black!55}#1\par}}
\newcommand{\gapcard}[2]{{\bodyfont\textbf{\textcolor{ICLBlue}{#1}}~— #2\par}\vspace{1.4mm}}

% Centre an image without the vertical padding of the center environment.
\newcommand{\posterfig}[1]{{\centering#1\par}}

%----------------------------------------------------------------------------------------
%	POSTER INFORMATION
%----------------------------------------------------------------------------------------

\postertitle{\raggedright Where to Measure, What to Measure With\\[6pt]
	{\ImperialSansMedium\fontsize{26pt}{30pt}\selectfont Evidence-Aware Drinking-Water Monitoring Under Uncertainty}}

\posterauthors{Muhammad Faseeh Memon\\[6pt]
	{\fontsize{20pt}{24pt}\selectfont MSc Environmental Engineering \quad Supervisor: Professor Michael Templeton\\
	Department of Civil and Environmental Engineering, Imperial College London}}

% No coauthor logos and no funder logos: this work has neither.

%----------------------------------------------------------------------------------------

\begin{document}

\setlength{\parskip}{0pt}

\titlesection

\vspace{0.5mm}

\noindent
%=========================================================================
%	LEFT KNOWLEDGE RAIL - 18.5 cm
%=========================================================================
\begin{minipage}[t]{18.5cm}\vspace{0pt}

	\railhead{Introduction — why this matters}

	\callout{363 million}
	\vspace{-1.5mm}
	{\bodyfont people in India lack safely managed drinking water.~[1]\par}

	\vspace{1.5mm}

	\callout{35 / 38 surveys}
	\vspace{-1.5mm}
	{\bodyfont found fewer households free of \textit{E.\ coli} at point of use than at collection; mean \textbf{−12.9 pp}.~[2]\par}

	\vspace{1.5mm}

	{\bodyfont\textbf{\textcolor{ICLDarkBlue}{Where to test, and what to test with?}}\par}

	\vspace{2mm}

	\railhead{Literature review — what was missing}

	{\bodyfont\textbf{\textcolor{ICLBlue}{DATA}} units, chemical basis and denominators differ across archives. \textbf{\textcolor{ICLBlue}{PREDICTION}} random validation overstates new-area performance. \textbf{\textcolor{ICLBlue}{TECHNOLOGY}} nominal capability is not exact-product evidence.\par}

	\vspace{1.5mm}
	{\bodyfont\textbf{\textcolor{ICLDarkBlue}{Gap:}} no study connected all three into a pathway permitted to stop and request evidence.~[3--5]\par}

	\vspace{2mm}

	\railhead{Methodology}

	\posterfig{\includegraphics[width=12.8cm, trim=5pt 100pt 5pt 5pt, clip]{poster_final/methodology_figure_3_1.pdf}}

	\vspace{1.6mm}
	{\capfont\textbf{Five-stage evidence-to-decision sequence;} abstention is a designed output. \textcolor{black!55}{Source: thesis Fig. 3.1.}\par}

	\vspace{1mm}

	\railhead{References}

	{\reffont\raggedright
	\begin{minipage}[t]{8.6cm}\raggedright\vspace{0pt}
	[1]~WHO \& UNICEF (2023). \textit{Progress on household drinking water, 2000--2022}.\par\vspace{1mm}
	[2]~Santos, T. \textit{et al.} (2023). \textit{E.\ coli} in drinking water sources. \textit{J.\ Water Health}.\par\vspace{1mm}
	[3]~Roberts, D. \textit{et al.} (2017). Cross-validation for structured data. \textit{Ecography}.\par\vspace{1mm}
	[4]~Camargo, E. \textit{et al.} (2023). Low-cost water sensors for IoT. \textit{Sensors}.\par
	\end{minipage}\hspace{0.7cm}%
	\begin{minipage}[t]{8.6cm}\raggedright\vspace{0pt}
	[5]~Pichel Mira, N. \textit{et al.} (2023). Field testing in LMIC regions. \textit{J.\ Environ.\ Chem.\ Eng.}\par\vspace{1mm}
	[6]~Trimmer, J. \textit{et al.} (2024). Costs of monitoring in rural Africa. \textit{Environ.\ Sci.\ Technol.}\par\vspace{1mm}
	[7]~Van Geen, A. \textit{et al.} (2019). Testing 30,000 wells for arsenic. \textit{Sci.\ Total Environ.}\par
	\end{minipage}
	}

\end{minipage}%
\hspace{1.0cm}%
%=========================================================================
%	MAIN RESULTS FIELD - 59.4 cm
%=========================================================================
\begin{minipage}[t]{59.4cm}\vspace{0pt}

	%--------------------------------------------------------------
	%	TOP ORIENTATION SPINE - P2 GRAPHICAL ABSTRACT
	%--------------------------------------------------------------
	\posterfig{\includegraphics[width=34.5cm]{poster_final/P2_graphical_abstract.png}}

	\vspace{1mm}
	{\srcfont\color{black!55} Thesis-derived synthesis; values verified against the submitted report.\par}

	\vspace{3mm}

	\mainhead{Results \& Discussion — what survived, what transferred, and where measurement matters next}

	\vspace{3mm}

	%--------------------------------------------------------------
	%	MAIN RESULTS ROW  (21.6 + 1.0 + 36.4 = 59.0 < 59.4)
	%--------------------------------------------------------------
	\noindent
	\begin{minipage}[t]{19.8cm}\vspace{0pt}
		\includegraphics[width=19.8cm]{poster_final/P1_decision_uncertainty_map.png}
		\srctag{State/UT-level measurement priority, not contamination severity or population prevalence. Source: Fig. 4.7 / Table 4.11.}
	\end{minipage}%
	\hspace{1.0cm}%
	\begin{minipage}[t]{38.2cm}\vspace{0pt}
		\blockhead{Honest validation: what survives geographic separation?}

		\posterfig{\includegraphics[width=23.0cm, trim=0pt 405pt 18pt 6pt, clip]{poster_final/result_figure_4_4.pdf}}

		\vspace{2mm}
		{\bodyfont\textbf{\textcolor{ICLDarkBlue}{5 / 7}} tasks retain acceptable-or-better spatial discrimination; label shuffle collapses to \textbf{\textcolor{ICLDarkBlue}{ROC-AUC 0.501}}. \textbf{Random splitting overstated deployment-facing performance; useful signal remained for most tasks.}\par}
		\srctag{Source: submitted thesis, Fig. 4.4 panel (a), Tables 4.3--4.4 and 4.7.}
	\end{minipage}

	\vspace{3mm}

	%--------------------------------------------------------------
	%	BOTTOM RESULTS ROW  (18.0 + 0.8 + 15.0 + 0.8 + 24.6 = 59.2 < 59.4)
	%--------------------------------------------------------------
	\noindent
	% --- Block 1: India transfer -----------------------------------
	\begin{minipage}[t]{18.0cm}\vspace{0pt}
		\blockhead{External India transfer: ranking survives where calibration does not}
		\includegraphics[width=18.0cm, trim=0pt 42pt 12pt 5pt, clip]{poster_final/result_figure_4_5.pdf}
		\srctag{Conductivity, TDS, fluoride and nitrate retain useful ordering; pH routes to direct measurement, arsenic to evidence acquisition. Local recalibration precedes probability-based use. Source: Fig. 4.5 / Table 4.6.}
	\end{minipage}%
	\hspace{0.8cm}%
	% --- Block 2: binding constraint (figure carries its own title) --
	\begin{minipage}[t]{15.0cm}\vspace{0pt}
		\includegraphics[width=15.0cm, trim=20pt 45pt 15pt 8pt, clip]{poster_final/result_figure_4_13.pdf}
		\srctag{Data constrain arsenic, turbidity and \textit{E.\ coli}; pH at transfer; the four rank-supported parameters at calibration and product evidence. Source: Fig. 4.13.}
	\end{minipage}%
	\hspace{0.8cm}%
	% --- Block 3: application + conclusion --------------------------
	\begin{minipage}[t]{24.6cm}\vspace{0pt}
		\blockhead{From analysis to monitoring action}

		{\bodyfont\textbf{\textcolor{ICLDarkBlue}{Rajasthan:}} conductivity 12,349 records, \textbf{54.06\%} above the 1,500\,µS/cm convention; TDS 3,131 records, \textbf{51.96\%} above 1,000\,mg/L.\par}

		\vspace{1.6mm}
		{\bodyfont\textbf{\textcolor{ICLBlue}{Rank → confirm → recalibrate locally → technology-class shortlist → product-level abstention}}\par}

		\vspace{1.5mm}
		{\ImperialSansExtraBold\fontsize{26pt}{29pt}\selectfont\color{ICLDarkBlue}Only 2 / 16 research priorities concern sensor design.\par}
		\vspace{1mm}
		{\capfont Most need evidence existing classes can supply: threshold-near validation • recalibration labels • cost and energy per valid result • reliability • operator effect • persistent identifiers.~[6,\,7]\par}

		\vspace{1.6mm}
		{\color{RuleGrey}\rule{\linewidth}{1.5pt}}\par
		\vspace{1.2mm}

		{\bodyfont
		\textbf{1.}~Honest validation preserves decision-useful signal.\par\vspace{0.8mm}
		\textbf{2.}~Rank transfer supports targeted confirmation, 4 of 6.\par\vspace{0.8mm}
		\textbf{3.}~All 26 contexts return a defined next action.\par}

		\vspace{1.8mm}
		{\ImperialSansSemiBold\fontsize{20pt}{23pt}\selectfont\color{ICLBlue}Prediction prioritises measurement; evidence sets how far the decision may go.\par}
	\end{minipage}

\end{minipage}

\end{document}
